%! The Statistical Cycle and Types of Data — Unit 1, Lesson 1.1 — Homework
% PILOT SHAPE (2026-09-06), Main Ideas / Notes table (2026-09-12).
% This is the lesson's SCORED (1 to 5) individual practice and the LAST component in the packet.
% Students START IT IN CLASS in the last ten minutes of the period, alone, while the
% teacher circulates for the second formative read; they finish it at home. Packet pages are
% the default; a Desmos activity may replace them on a given day (decided at Close & Assign).
% THIRD CONTEXT: the notes teach on the Riverbend Student Council's bike-rack survey, the
% Guided Practice runs on the council's SECOND survey, and these items are on the Lakeside
% Farmers Market — recall is not the skill. Lakeside is also Lesson 1.2's teaching context,
% so the set doubles as the hand-off.
% TWO PAGES MAX (user decision 2026-09-06): two boxes — items 1-2 on page 1, items 3-6 and
% the spiralbox on page 2 — so no item breaks across a page.
% Item 4 rows (i) and (ii) are the deliberate CONTRAST PAIR: the SAME shoppers, two questions,
% two types — the target misconception (that the data, not the question, decide the type).
% Item 6 is the crux head-on: every number in the sentence is correct and the conclusion
% still does not follow, and the problem entered at stage 1.
% THE DATA — Lakeside Farmers Market: about 800 Saturday shoppers; 50 asked on one Saturday ->
%   26 drove (52%), 12 walked (24%), 7 biked (14%), 5 bus (10%); 0.10 x 800 = 80.
% PAGE LOCKSTEP: the key is generated from this file by templates/lesson/mkkey.py (spec:
% templates/lesson/examples/lesson01_homework_key.json); every work block is byte-identical.
\documentclass[12pt]{article}
\usepackage{saar-article}
\usepackage{saar-boxes}
\usepackage{grading}
\gradingassignment{saar-u01-l01-hw}
\gradingcourse{stats_analysis_algebraic_reasoning}
\newcommand{\TallMath}[1]{$\displaystyle #1\rule[-1.4em]{0pt}{3.2em}$}
% Word bank strip for every fill-in cluster (user request 2026-09-06). Defined per document;
% the key inherits it and prints the same strip, so heights match.
\newcommand{\wordbank}[1]{\par\vspace{2pt}\noindent\fcolorbox{goldacc}{goldbg}{\parbox{\dimexpr\linewidth-2\fboxsep-2\fboxrule\relax}{\small\textbf{Word bank:} #1}}\par\vspace{4pt}}

\begin{document}

\pageheader{Unit 1, Lesson 1.1}{Homework\strut}
% NAMESTRIP: no \namedateperiod here. The student writes their name once, on the
% cover this component is stapled behind.

\vspace{0.12in}

% Authored identically in homework/ and homework_key/.
\begin{remindbox}
\small
\textbf{This is your graded homework.} It is scored 1 to 5 and due at the start of the
first class after two study halls. You will start it in the last minutes of the period, so ask about anything that stalls
you before you leave, then finish it on your own.
\end{remindbox}

\vspace{0.08in}

\boxguard[20]
\begin{scenariobox}[Lakeside Farmers Market]{cerulean}
About $800$ shoppers come to the Lakeside Farmers Market on a Saturday morning. The manager is
deciding whether to ask the town for a free shuttle from the library parking lot, so she asks
$50$ shoppers one Saturday: \emph{``How did you travel to the market today?''}

\vspace{2pt}
Her $50$ answers: $26$ drove, $12$ walked, $7$ biked, $5$ rode the bus.
\end{scenariobox}

\vspace{0.08in}

\boxguard
\begin{notesbox}{The Data Cycle at the Market}
Every answer names the group it is about --- the $50$ shoppers she asked, or all $800$.
\wordbank{categorical \;\textbullet\; quantitative \;\textbullet\; how the shopper traveled \;\textbullet\; 50 \;\textbullet\; 52\% \;\textbullet\; 24\% \;\textbullet\; 14\% \;\textbullet\; 10\% \;\textbullet\; 100\%}
\begin{enumerate}[leftmargin=1.6em, itemsep=9pt, topsep=3pt]

  \item Name the \textbf{variable of interest} --- what she records for each shopper --- and
        its type.

        \par\vspace{4pt}
        Variable: \blankbox{\gradedblank[id=q1a, type=short, points=0.5, answer={how the shopper traveled}, check=text]{6cm}{13.5pt}} \hfill Type: \blankbox{\gradedblank[id=q1b, type=short, points=0.5, answer={categorical}, check=text]{3.4cm}{13.5pt}}
        \par\vspace{8pt}
        \gradedresponse[How do you know?]{\gradedblank[id=q1c, type=rubric, points=1, display={the answers are labels --- no arithmetic on them means anything}]{\linewidth}{2in}}

  \item Complete her frequency table, then show your work for the \textbf{bus} row.

        \begin{center}
        \renewcommand{\arraystretch}{1.45}
        \begin{tabularx}{\linewidth}{@{} l c X @{}}
        \toprule
        \textbf{How they traveled} & \textbf{Count} & \textbf{Relative frequency (percent)} \\
        \midrule
        Drove  & 26 & \blank{4.2cm}\hspace{0.2cm}\smash{\gradedblank[id=q2a, type=short, points=0.25, answer={52}, check=numeric, display={52\%}]{2cm}{18pt}} \\
        Walked & 12 & \blank{4.2cm}\hspace{0.2cm}\smash{\gradedblank[id=q2b, type=short, points=0.25, answer={24}, check=numeric, display={24\%}]{2cm}{18pt}} \\
        Biked  &  7 & \blank{4.2cm}\hspace{0.2cm}\smash{\gradedblank[id=q2c, type=short, points=0.25, answer={14}, check=numeric, display={14\%}]{2cm}{18pt}} \\
        Bus    &  5 & \blank{4.2cm}\hspace{0.2cm}\smash{\gradedblank[id=q2d, type=short, points=0.25, answer={10}, check=numeric, display={10\%}]{2cm}{18pt}} \\
        \midrule
        \textbf{Total} & \blankbox{\gradedblank[id=q2e, type=short, points=0.5, answer={50}, check=numeric]{1.4cm}{18pt}} & \blank{4.2cm}\hspace{0.2cm}\smash{\gradedblank[id=q2f, type=short, points=0.5, answer={100}, check=numeric, display={100\%}]{2cm}{18pt}} \\
        \bottomrule
        \end{tabularx}
        \end{center}

        \begin{work}
          \dfrac{5}{50} &= 0.10 \\
                        &= 10\%
        \end{work}

\end{enumerate}
\end{notesbox}

\boxguard[30]
\begin{notesbox}{The Data Cycle at the Market, continued}
\begin{enumerate}[leftmargin=1.6em, itemsep=5pt, topsep=2pt, start=3]

  \item Estimate how many of the market's $800$ Saturday shoppers ride the bus.

        \begin{work}
          0.10 \times 800 &= 80 \text{ shoppers}
        \end{work}

        Write that result as one sentence for the manager's report: name the group she asked,
        the number, and what was measured.
        \gradedresponse{\gradedblank[id=q3, type=rubric, points=2, display={Of the 50 shoppers asked one Saturday, 10\% travel to the market by bus. If they are typical, that is about 80 of the market's 800 Saturday shoppers.}]{\linewidth}{2in}}


  \item The manager could ask her $50$ shoppers either question below. \textbf{Decide the type
        from the question alone} --- nothing has been collected yet.

        \begin{center}
        \renewcommand{\arraystretch}{1.3}
        \begin{tabularx}{\linewidth}{@{} X p{2.6cm} @{}}
        \toprule
        \textbf{The question} & \textbf{Type} \\
        \midrule
        (i) ``How did you travel to the market today?''
          & \blankbox{\gradedblank[id=q4a, type=short, points=0.25, answer={categorical}, check=text]{2.2cm}{17pt}} \\
        (ii) ``How many miles did you travel to the market today?''
          & \blankbox{\gradedblank[id=q4c, type=short, points=0.25, answer={quantitative}, check=text]{2.2cm}{17pt}} \\
        \bottomrule
        \end{tabularx}
        \end{center}

        How would you summarize the answers to each question?
        \gradedresponse[\textbf{(i)} Summarize with\dots]{\gradedblank[id=q4b, type=rubric, points=0.25, display={counts and percents}]{\linewidth}{2in}}
        \gradedresponse[\textbf{(ii)} Summarize with\dots]{\gradedblank[id=q4d, type=rubric, points=0.25, display={a mean}]{\linewidth}{2in}}

        Same $50$ shoppers, same Saturday. Explain what made the two types different.
        \gradedresponse{\gradedblank[id=q4e, type=rubric, points=1, display={The question did --- one records a label, the other records a number of miles.}]{\linewidth}{2in}}

  \item Each mistake below broke the study at one stage. Write the number of the stage the
        problem \textbf{entered} at --- not the stage where the bad sentence showed up.

        \begin{enumerate}[leftmargin=1.6em, itemsep=4pt, topsep=3pt, label=(\alph*)]
          \item The original question never said \emph{which} shoppers or \emph{which} day.
                \hfill \blankbox{\gradedblank[id=q5a, type=short, points=0.5, answer={1}, check=numeric]{1.2cm}{13.5pt}}
          \item She asked only shoppers walking in from the parking lot. \hfill \blankbox{\gradedblank[id=q5b, type=short, points=0.5, answer={2}, check=numeric]{1.2cm}{13.5pt}}
          \item The $50$ answers were never counted by travel method. \hfill \blankbox{\gradedblank[id=q5c, type=short, points=0.5, answer={3}, check=numeric]{1.2cm}{13.5pt}}
          \item The write-up reports $70\%$ drove when the table says $52\%$.
                \hfill \blankbox{\gradedblank[id=q5d, type=short, points=0.5, answer={4}, check=numeric]{1.2cm}{13.5pt}}
        \end{enumerate}

  \item Her report says: \textit{``Only $10\%$ of our shoppers ride the bus, so a free shuttle
        would be used by very few people.''} \textbf{Every number in that sentence is
        correct.} Explain why the conclusion still does not follow, and name the stage the
        problem entered at.
        \gradedresponse{\gradedblank[id=q6, type=rubric, points=2, display={Nobody was ever asked about a shuttle, so the survey cannot say who would use one; 10\% of 800 is still about 80 shoppers. It entered at stage 1 --- the wrong question.}]{\linewidth}{2in}}

\end{enumerate}
\end{notesbox}

\boxguard[3]   % the two-line spiralbox needs about three baselines; a larger guard pushes it to a third page
\begin{spiralbox}
\textbf{Next lesson: Populations, Samples, Parameters, and Statistics.} She asked $50$ shoppers
and then wrote about all $800$. Both groups have names --- and so do the numbers from each.
\end{spiralbox}

\end{document}
